\chapter{振動子の波動関数、スピン1/2、粒子交換}

\section{至る所に振動子}

安定な平衡をもつ系は、その平衡から変位させられると振動しうる。振動数は一つかもしれないし、多数かもしれない。たとえばバイオリン弦は基本モードと高次倍音をもつ。弦を弾くと、それらのモードのある組み合わせが励起される。

ここでより重要な例は場である。電磁場の平衡配置は真空である。すなわち電場も磁場もなく、電磁振動もない。空洞や導波管に放射を送り込むと、場は許されたモードで振動しうる。各モードはそれぞれ固有の振動数をもつ。量子論では、それらのモードは光子の集まりを記述する。

同じ考え方は結晶にも当てはまる。すべての原子を独立な振動子と描くのは誤解を招く。原子は互いに力を及ぼし合うので、一つの原子の変位は隣の原子に影響し、波が結晶中を伝わる。その結合した運動の基準モードは調和振動子であり、その量子はフォノンと呼ばれる。光子とフォノンは異なる物理的対象であるが、どちらも同じ振動子代数に立ち返る。

\section{振動子代数}

振動子の議論では、質量と換算プランク定数を1に等しく置く。
\[
m=1,\qquad \hbar=1.
\]
ハミルトニアンは
\begin{equation}
H=\frac{p^{2}}{2}+\frac{\omega^{2}x^{2}}{2}.
\end{equation}
である。パラメータ\(\omega\)はポテンシャルの曲率を定め、同時に振動数でもある。

生成消滅演算子を
\begin{equation}
A_{\pm}=\frac{p\pm i\omega x}{\sqrt{2\omega}}.
\end{equation}
で定義する。規格化は
\begin{equation}
[A_{-},A_{+}]=1.
\end{equation}
となるように選ばれている。数演算子は
\begin{equation}
N=A_{+}A_{-}.
\end{equation}
である。これらの式を合わせると、\(A_{+}\)は\(N\)の固有値を1だけ上げ、\(A_{-}\)はそれを1だけ下げることがわかる。得られるスペクトルは
\begin{equation}
n=0,1,2,3,\ldots .
\end{equation}
である。

二つの梯子演算子を掛け合わせると、古典的エネルギーに対する量子補正も現れる。
\begin{align}
N
&=\frac{1}{2\omega}(p+i\omega x)(p-i\omega x)\\
&=\frac{1}{2\omega}\left(p^{2}+\omega^{2}x^{2}
 +i\omega[x,p]\right)\\
&=\frac{1}{2\omega}\left(p^{2}+\omega^{2}x^{2}-\omega\right),
\end{align}
ここで\([x,p]=i\)である。したがって、
\begin{equation}
H=\omega\left(N+\frac12\right).
\end{equation}
である。付加的な定数が問題にとって無関係なら、シフトしたハミルトニアン
\begin{equation}
H_{\mathrm{shifted}}=H-\frac{\omega}{2}=\omega N.
\end{equation}
を用いるのが便利である。項\(\omega/2\)は振動子の零点エネルギー、すなわち基底状態エネルギーである。場のモードでは、対応する項は真空エネルギーと呼ばれる。これは不確定性原理を反映している。すなわち、\(x\)と\(p\)の両方が同時に鋭くゼロに等しい状態は存在しえないので、\(H\)における二つの正の二次項がともにゼロの期待値をもつことはできない。

\begin{classroomqa}
\lecturetimestamp{00:09:46}
\audiencequestion{なぜゼロが数演算子のスペクトルに含まれるのか。}
\lecturerresponse{それは基底状態の波動関数を構成することによって明示的に示すことができる。その構成は\(N=0\)をもつ状態を提示するであろう。}
\end{classroomqa}

その構成を始めるために、状態ベクトルを位置の波動関数に置き換える。エネルギー固有状態は
\[
\psi_{0}(x),\ \psi_{1}(x),\ \psi_{2}(x),\ldots,
\]
となり、基本演算子は
\begin{equation}
X\longmapsto x,\qquad
P\longmapsto -i\frac{d}{dx}.
\end{equation}
として作用する。\(\hbar\)を復元すれば、第二の置き換えは
\(P\mapsto-i\hbar\,d/dx\)
となるであろう。

\(\psi_{0}\)を解く前に、この問いにはより速い抽象的答えがある。規格化可能な振動子固有状態の離散的な梯子の上で、\(N=A_{+}A_{-}\)は非負であるから、単位刻みでの繰り返しの下降は無限に続くことはできない。端点---その規格化可能な波動関数は以下で構成される---には、\(\lvert0\rangle\)と記される状態が存在しなければならず、それに対して
\begin{equation}
A_{-}\lvert0\rangle=\mathbf 0.
\end{equation}
が成り立つ。そうでなければ、\(A_{-}\lvert0\rangle\)がもう一つ一段下の状態となるであろう。ただちに
\begin{equation}
N\lvert0\rangle
=A_{+}A_{-}\lvert0\rangle
=\mathbf 0
=0\,\lvert0\rangle.
\end{equation}
が従う。したがって、最下段の状態は数固有値ゼロをもつ。

\begin{classroomqa}
\lecturetimestamp{00:14:57}
\audiencequestion{ゼロが固有値なら、右辺はゼロ固有ベクトルを含むべきか。}
\lecturerresponse{ケット\(\lvert0\rangle\)は規格化された非零ベクトルであり、そのラベルはゼロである。\(N\lvert0\rangle=0\,\lvert0\rangle\)において、右辺はゼロベクトルである。基底状態のケットとゼロベクトルは異なる対象である。}
\end{classroomqa}

\section{Schr\"odinger方程式と物理的解}

ここで位置表示に戻る。時間に依存しないSchr\"odinger方程式は
\begin{equation}
\left[
-\frac12\frac{d^{2}}{dx^{2}}
+\frac12\omega^{2}x^{2}
\right]\psi(x)=E\psi(x).
\end{equation}
であり、二階微分方程式として、\(E\)の任意の値に対して形式的な解をもつ。それらの形式的な解のほとんどは物理的な振動子状態ではない。それらは大きな\(\lvert x\rvert\)で指数関数的に増大するため、そのノルムは発散する。

束縛状態に対しては、波動関数は二乗可積分でなければならない。
\begin{equation}
\int_{-\infty}^{\infty}\psi^{*}(x)\psi(x)\,dx<\infty.
\end{equation}
そのとき、積分が1に等しくなるように規格化できる。そのような解を生じるのは特別なエネルギー値だけである。無限直線上の位置状態と運動量状態はこの点で例外的であるが、指数関数的に増大する振動子解は単に状態空間から排除される。

\section{ガウス型基底状態}

二階の方程式も解けるが、はしご代数を使えばより簡単な一階の方程式が得られる：
\begin{equation}
A_{-}\psi_{0}(x)=0.
\end{equation}
右辺がゼロのとき、\(A_{-}\)の全体の因子は無関係なので、これは次と等価である：
\begin{equation}
(p-i\omega x)\psi_{0}(x)=0.
\end{equation}
\(p=-i\,d/dx\)を用いると
\begin{equation}
\left(\frac{d}{dx}+\omega x\right)\psi_{0}(x)=0.
\end{equation}

\begin{classroomqa}
\lecturetimestamp{00:22:32}
\audiencequestion{これは変数分離で解けるのだろうか？}
\lecturerresponse{独立変数が一つしかないので、分離すべき二つの変数は存在しない。\(\psi_{0}'/\psi_{0}=-\omega x\)と書けば、この方程式は対数微分の形をしている。そのため波動関数を指数関数として書くことが示唆される。}
\end{classroomqa}

これを解く有用な方法は次のように書くことである：
\begin{equation}
\psi_{0}(x)=e^{f(x)}.
\end{equation}
すると
\[
\psi_{0}'(x)=f'(x)e^{f(x)},
\]
となり、共通の指数因子が相殺される。方程式は次に帰着する：
\begin{equation}
\frac{df}{dx}+\omega x=0.
\end{equation}

\begin{classroomqa}
\lecturetimestamp{00:26:28}
\audiencequestion{\(f(x)\)にマイナス符号が付くべきだろうか？}
\lecturerresponse{はい。\(df/dx=-\omega x\)なので、積分すると負の二次項が得られる。}
\end{classroomqa}

したがって
\begin{equation}
f(x)=-\frac12\omega x^{2}+C,
\end{equation}
であり、
\begin{equation}
\psi_{0}(x)\propto e^{-\omega x^{2}/2}.
\end{equation}
である。乗法的な規格化定数は指定しないままにしておく。重要な点はガウス型の形状である：滑らかで、原点付近に集中し、大きな\(\lvert x\rvert\)では急速に減少する。明らかに二乗可積分である。

そのエネルギーを知るために答えを二階の方程式に代入し直す必要はない。代数がすでに
\begin{equation}
E_{0}=\frac{\omega}{2}.
\end{equation}
を与えている。

\section{励起状態、節、古典的運動}

第一励起状態は基底状態を上昇させて得られる：
\begin{equation}
\psi_{1}(x)\propto A_{+}\psi_{0}(x)
\propto(p+i\omega x)e^{-\omega x^{2}/2}.
\end{equation}
基底状態の方程式は
\[
p\psi_{0}=i\omega x\,\psi_{0}.
\]
と言っている。はしご演算子の符号を変えると、相殺が和に変わる：
\begin{equation}
(p+i\omega x)\psi_{0}
=2i\omega x\,\psi_{0}.
\end{equation}
全体の定数を除いて、
\begin{equation}
\psi_{1}(x)\propto x\,e^{-\omega x^{2}/2}.
\end{equation}

この波動関数は奇である：
\begin{equation}
\psi_{1}(-x)=-\psi_{1}(x).
\end{equation}
その確率密度は偶であるが、原点に節を持つ：
\begin{equation}
\lvert\psi_{1}(0)\rvert^{2}=0.
\end{equation}
したがって第一励起状態の確率密度は\(x=0\)に節を持つ。

\(A_{+}\)を繰り返し作用させると、同じガウス型の包絡線に掛かるより高次の多項式が得られる。\footnote{講義では、より高次の多項式因子、各段階で追加される節、そして\(00{:}34{:}25\)から\(00{:}35{:}49\)への対称性の交代が述べられている。ガウス型掛ける多項式という形は、Leonard SusskindとArt Friedmanの\emph{Quantum Mechanics: The Theoretical Minimum}（2014）の振動子固有関数の議論によって裏付けられている。}規格化を除いて、
\begin{equation}
\psi_{n}(x)=P_{n}(x)e^{-\omega x^{2}/2}.
\end{equation}
パリティは偶と奇の間で交代し、各段階で節が一つ加わる。\(n\)が増大するにつれて、波動関数は原点からより遠くへ広がり、より速い振動を発達させる。

この二つの特徴には直接的な物理的意味がある。原点からより遠くへの広がりは、ポテンシャルが\(x^{2}\)として増大するため、より大きなポテンシャルエネルギーに対応する。位置空間での速い振動はより大きな運動量を示す。高度に励起された振動子は中心からより遠くへ到達する可能性が高く、中心領域を通過するときにはより大きな運動量を運び、その波動関数はそこではより速く振動する。

\begin{classroomqa}
\lecturetimestamp{00:37:18}
\audiencequestion{\(n\)が大きくなっても調和振動子は古典的振る舞いに近づかないのだろうか？}
\lecturerresponse{古典的運動は一つの定常エネルギー固有状態によって表されるものではない。それは多くのエネルギー準位を含む波束の中に現れる。たとえば変位したガウス型は古典的振動子のように往復して発展し、振れ幅が大きくなるにつれて相対的不確定性は小さくなる。}
\end{classroomqa}

基底状態のガウス関数から始め、それを原点から変位させる。変位した波束はもはやエネルギー固有状態ではなく、多くの固有状態の重ね合わせである。時間に依存するSchr\"odinger方程式の下で、その成分は異なる位相を獲得する。波束はポテンシャルの一方の側から他方の側へと揺れ動き、中心付近で急速な位相変化を生じ、転回点付近では再び滑らかになる。その振幅が波束の幅と比べて大きいほど、運動はより古典的に見える。

\begin{classroomqa}
\lecturetimestamp{00:39:29}
\audiencequestion{第二励起状態の波動関数は原点でゼロを通るべきだろうか？}
\lecturerresponse{いいえ。その多項式因子は単に\(x^{2}\)ではなく、定数項も含むので、波動関数は\(x=0\)で必ずしも消える必要はない。}
\end{classroomqa}

はしごの議論はまた、\(A_{+}\)と\(A_{-}\)が二乗可積分な振動子の定義域上でエルミート共役であり、随伴関係に必要な境界挙動を持つことを仮定している。指数関数的に増大する形式的解はその定義域に属さない。

\begin{classroomqa}
\lecturetimestamp{00:40:45}
\audiencequestion{ハミルトニアンがこれらの代数的方法で解けるかどうかを認識する一般的な方法はあるか。}
\lecturerresponse{そのような一般的な規則はない。厳密に解ける例は稀である。調和振動子と理想的な非相対論的水素原子は、代数的方法がうまくいく幸運な場合である。}
\end{classroomqa}

\section{角運動量としてのスピン1/2}

スピンは、粒子がその重心の軌道運動とは独立に持つ角運動量である。静止している粒子は軌道角運動量\(\mathbf R\times\mathbf P\)を持たないが、それでもスピンを持つことができる。小さな回転する物体を思い描こうとするか、スピンを内在的な量子特性として扱うことができる。決定的な記述は数学的である。その成分は空間回転の下で変換し、角運動量の交換関係を満たす。

この時点から\(\hbar\)を復元する。角運動量の成分は
\begin{equation}
[L_{z},L_{x}]=i\hbar L_{y},
\end{equation}
を満たし、他の2つの関係は巡回置換によって得られる。無次元のパウリ行列は
\begin{equation}
\begin{aligned}
\sigma_{z}&=
\begin{pmatrix}
1&0\\
0&-1
\end{pmatrix},\\[0.4em]
\sigma_{x}&=
\begin{pmatrix}
0&1\\
1&0
\end{pmatrix},\\[0.4em]
\sigma_{y}&=
\begin{pmatrix}
0&-i\\
i&0
\end{pmatrix}.
\end{aligned}
\end{equation}
である。この基底では、\(\sigma_{z}\)は対角的である。その固有値は\(+1\)と
\(-1\)である。

\begin{classroomqa}
\lecturetimestamp{00:47:35}
\audiencequestion{\(\sigma_x\)に負号はあるか。}
\lecturerresponse{いいえ。\(\sigma_x\)の2つの非対角成分はどちらも正である。}
\end{classroomqa}

\(00{:}49{:}35\)での板書計算は成分ごとの乗算を始める。積を完成させると
\begin{align}
\sigma_{z}\sigma_{x}
&=
\begin{pmatrix}
0&1\\
-1&0
\end{pmatrix}
=i\sigma_{y},\\
\sigma_{x}\sigma_{z}
&=
\begin{pmatrix}
0&-1\\
1&0
\end{pmatrix}
=-i\sigma_{y}.
\end{align}
したがって
\begin{equation}
[\sigma_{z},\sigma_{x}]=2i\sigma_{y}.
\end{equation}
余分な因子2は、物理的なスピン演算子を\footnote{と定義することによって除かれる。板書計算では\(\hbar\)が省略されている。ここでは\(S_i\)が角運動量の単位を持つように復元されている。}
\begin{equation}
S_{i}=\frac{\hbar}{2}\sigma_{i}.
\end{equation}
すると
\begin{equation}
[S_{z},S_{x}]=i\hbar S_{y},
\end{equation}
であり、巡回置換についても同様である。したがって、\(\hbar\)の因子を復元したパウリ行列の半分は、角運動量系を形成する。

\(\hbar=1\)と置くと、以下に示す形になる。\footnote{ソースフレームは\texttt{lecture\_04\_frame\_02.png}である。}
\begin{figure}[htbp]
\centering
\includegraphics[width=0.9\linewidth]{lecture_04_frame_02.png}
\caption{\(\hbar=1\)の単位では、パウリ行列の半分はスピン1/2の角運動量演算子であり、\([S_z,S_x]=iS_y\)を満たす。}
\lectureframe{00:55:41}
\end{figure}

2つの\(S_z\)固有値方程式は
\begin{equation}
S_z\lvert\pm\rangle
=\pm\frac{\hbar}{2}\lvert\pm\rangle.
\end{equation}
である。等価的に、無次元の磁気量子数は
\begin{equation}
m_{s}=\pm\frac12,
\qquad
S_{z}\lvert m_{s}\rangle
=m_{s}\hbar\,\lvert m_{s}\rangle.
\end{equation}
である。角運動量の昇降演算子は\(m\)を整数ステップでずらすが、これはすべてのはしごがゼロを通ることを要求するものではない。整数の多重項は整数の\(m\)値を含み、半整数の多重項は半整数値を含む。スピン1/2の場合、はしごは\(m_s=+1/2\)と\(m_s=-1/2\)の2つの要素だけを持つ。

これらの2つの状態はスピン1/2の二重項を形成する。電子、ミューオン、ニュートリノ、クォークは、そのような半整数スピン多重項で記述される粒子の中に含まれる。

\section{軌道角運動量とスピン角運動量}

ここで2種類の角運動量がある。軌道角運動量は原点に対する運動を記述し、
\begin{equation}
\mathbf L=\mathbf R\times\mathbf P,
\end{equation}
一方、\(\mathbf S\)は内在的スピンである。それらのベクトル和は全角運動量であり、
\begin{equation}
\mathbf J=\mathbf L+\mathbf S.
\end{equation}

\begin{classroomqa}
\lecturetimestamp{00:56:49}
\audiencequestion{系が全角運動量\(\mathbf J\)を持つとき、そのスピン\(\mathbf S\)はどのように同定されるか。}
\lecturerresponse{スピン量子数と軌道量子数は通常別々に与えられる。もし\(L=0\)ならば、すべての角運動量はスピンである。そうでなければ角運動量を加えるための量子規則が必要であり、それはここでは延期される。}
\end{classroomqa}

加算規則は、軌道角運動量とスピン、2つのスピン、または2つの軌道角運動量に等しく適用される。結合系はそれ自身の許された全角運動量の値を持つが、それらの規則はまだ必要ではない。

\begin{classroomqa}
\lecturetimestamp{00:58:18}
\audiencequestion{スピンの導入につながった実験的証拠は何か。}
\lecturerresponse{分光学が証拠を提供したが、周期表は特に透明な動機を与える。その殻構造は、電子が追加の2値的性質を持つことを要求する。}
\end{classroomqa}

\section{水素様縮退と殻の数え上げ}

軌道角運動量について、\(L^{2}\)固有値方程式は
\begin{equation}
L^{2}\lvert\ell,m\rangle
=\hbar^{2}\ell(\ell+1)\lvert\ell,m\rangle,
\end{equation}
であり、\(\ell\)の固定された値は
\begin{equation}
2\ell+1
\end{equation}
個の磁気状態を持ち、\(m=-\ell,-\ell+1,\ldots,\ell\)に対応する。

理想的な非相対論的クーロン問題には追加の縮退がある。その束縛状態のエネルギーは負であり、下方からゼロに近づくが、計数のためには準位の並びだけが重要である。最初の\(\ell=1\)準位は次の\(\ell=0\)準位と並び、最初の\(\ell=2\)準位は次の主殻に加わり、このパターンが続く。\footnote{\texttt{lecture\_04\_frame\_03.png}、\(01{:}02{:}01\)。完成した主殻の並びは\(01{:}02{:}33\)から\(01{:}04{:}10\)まで述べられている。}

\begin{figure}[htbp]
\centering
\begin{tikzpicture}[x=1.55cm,y=1.0cm]
\draw[->] (-0.65,0) -- (-0.65,4.35);
\draw[->] (-0.85,0) -- (4.75,0);
\node at (-0.95,4.08) {\(E\)};
\node at (4.55,-0.28) {\(\ell\)};
\node at (0,-0.30) {\(0\)};
\node at (2,-0.30) {\(1\)};
\node at (4,-0.30) {\(2\)};

\draw (-0.20,0.75) -- (0.20,0.75);

\draw (-0.20,2.00) -- (0.20,2.00);
\foreach \dx in {-0.18,0,0.18}
  \draw (2+\dx-0.06,2.00) -- (2+\dx+0.06,2.00);

\draw (-0.20,3.15) -- (0.20,3.15);
\foreach \dx in {-0.18,0,0.18}
  \draw (2+\dx-0.06,3.15) -- (2+\dx+0.06,3.15);
\foreach \dx in {-0.32,-0.16,0,0.16,0.32}
  \draw (4+\dx-0.05,3.15) -- (4+\dx+0.05,3.15);

\node[anchor=west] at (0.35,0.75) {\(n=1\)};
\node[anchor=west] at (2.35,2.00) {\(n=2\)};
\node[anchor=west] at (4.40,3.15) {\(n=3\)};
\end{tikzpicture}
\caption{最初の三つの理想的なクーロン殻は、等しい主エネルギーをもつ軌道多重項を組み合わせる。すなわち、第一殻に一つの状態、第二殻に\(1+3\)、第三殻に\(1+3+5\)である。}
\lectureframe{01:02:01}
\end{figure}

主殻\(n\)において、許される軌道の値は\footnote{このコンパクトな\(\ell\)範囲と以下の和が、\(01{:}02{:}33\)から\(01{:}04{:}10\)まで述べられた準位の並びと計数\(1,4,9,16,\ldots\)を再構成する。スペクトル図の始まりは\texttt{lecture\_04\_frame\_03.png}の\(01{:}02{:}01\)に現れる。}
\[
\ell=0,1,\ldots,n-1.
\]
したがって空間状態の数は
\begin{equation}
\sum_{\ell=0}^{n-1}(2\ell+1)=n^{2},
\end{equation}
となり、数列
\begin{equation}
1,\ 4,\ 9,\ 16,\ldots .
\end{equation}
を与える。これは軌道の計数のみである。各空間状態は二つのスピン値を許すので、対応する理想化された総数は
\begin{equation}
2,\ 8,\ 18,\ 32,\ldots=2n^{2}.
\end{equation}

\section{排他と二値的電子性質}

ヘリウムを独立電子近似で始める。その原子核は陽子電荷の二倍をもち、それによって水素様準位が再スケールされ、電子が内側へ引き寄せられる。電子--電子相互作用を今は無視すると、二つの電子が最低空間軌道を占めることができ、ヘリウム基底状態の妥当な第一近似を与える。

なぜそこに第三の電子も置かないのか。そうすればその軌道エネルギーは下がるが、実際には起こらない。制限はパウリの排他原理である。すなわち、二つの電子が同じ完全な一粒子状態を占めることはできない。これを通常の斥力と呼ぶのは誤解を招く。それは許される量子状態に対する制限である。

見かけの矛盾がある。二つのヘリウム電子は確かに同じ空間軌道を共有している。解決は、完全な電子状態が軌道波動関数以上のものを含むことにある。それはまた二値的な自由度を含む。その値を上と下と呼ぶ。最低軌道の一方の電子はスピン上、他方はスピン下でありうる。それらの空間状態は一致するが、完全な量子状態は一致しない。

この時点で計数はリチウムに向かって進み、次に引き返す。この引き返しは、なぜ排他とスピンの両方が必要かを明らかにする。スピンが全くない排他規則を想像せよ。最低空間状態を占められる電子は一つだけである。第二の電子は次の主殻の四つの空間状態の一つに入らねばならず、観測されるヘリウム基底状態が一意であるところにいくつかの代替を生じる。そのスピンなし版の排他は誤りである。

二つのスピン値を復活させる。二つの電子が最低軌道を反対向きのスピンで満たす。第三の電子は、中性リチウムのように、次の主殻に入らねばならない。その殻には一つの\(\ell=0\)空間状態と三つの\(\ell=1\)空間状態がある。スピンを数える前には四つの可能性がある。

\begin{classroomqa}
\lecturetimestamp{01:14:58}
\audiencequestion{どのようにして\(s\)状態と\(p\)状態が第一励起殻で八つの可能性を与えるのか。}
\lecturerresponse{\(\ell=0\)をもつ空間状態が一つ、\(\ell=1\)をもつものが三つあり、四つの空間状態を与える。それぞれが二つのスピン値を許すので、総数は\(2(1+3)=8\)である。}
\end{classroomqa}

これが周期表の第二行の八つの項目の背後にある単純な計数である。これは意図的に理想化であり、後の議論は電子--電子相互作用がすぐに独立水素様殻の描像を損なうことを明らかにするであろう。

二値的性質は角運動量としても検証できる。電子のスピンには磁気モーメントが伴うので、磁場はその二つのスピン射影に異なるエネルギーを与える。ゼロ磁場で縮退していた準位は二つに分裂し、小さな磁気モーメントが磁場に平行または反平行に整列する場合に予想される通りになる。この分裂は、加えられた二値的自由度が任意の二値ラベルではなく角運動量として振る舞うかどうかを検証する。

\section{光子、電子、そして殻描像の限界}

粒子がすべて同じ占有規則に従うわけではない。多数の光子が同一の電磁モードを占めてよい。古典的な光波は、そのモードに多数の光子を含む状態として理解してよい。電子は異なる振る舞いをする。排他によって、電子は一つの完全な量子状態に詰め込まれることが妨げられる。この対比が、二つの粒子統計の区別を動機づける。

\begin{classroomqa}
\lecturetimestamp{01:20:28}
\audiencequestion{原子に束縛された電子ではなく、自由な電子の流れの場合はどうなるのか。}
\lecturerresponse{排他規則は依然として適用される。電子はすべてが一つの量子状態を占めることはできない。レーザーは多数の光子を同一モードに置くことができるが、電子の流れをそのように準備することはできない。}
\end{classroomqa}

この制限は電子の性質であり、原子束縛に固有の特異性ではない。

\begin{classroomqa}
\lecturetimestamp{01:21:12}
\audiencequestion{軌道角運動量とは、特に原子核のまわりの回転のことか。}
\lecturerresponse{角運動量は、選ばれた原点に対して定義される。原子では原子核が自然な原点であるが、自由粒子は、どの原点を選んでもよいかに対して角運動量 \(\mathbf R\times\mathbf P\) をもつ。}
\end{classroomqa}

回転の議論はすべて、その回転がどの点のまわりに行われるかを選ぶ。原子は自然な中心を与える。その原子核が、ほぼ中心的なポテンシャルの中心だからである。自由空間ではどの原点も特別ではないが、それでも任意に選んだ原点のまわりで角運動量を計算することはできる。

\begin{classroomqa}
\lecturetimestamp{01:22:34}
\audiencequestion{周期表にはなぜ、8個の元素からなる行が繰り返された後、18個の元素からなる行が繰り返されるのか。}
\lecturerresponse{単純な独立電子の殻模型では、この全体的なパターンを説明できない。電子間の相互作用を無視しているため、この描像はすぐに破綻する。}
\end{classroomqa}


大きな原子では、外側の電子は原子核だけでなく、かなりの内殻電子雲とも相互作用する。すべての電子が遮蔽されていないクーロン場の中を独立に運動しているかのように扱うのは、もはや妥当ではない。

\begin{classroomqa}
\lecturetimestamp{01:23:33}
\audiencequestion{なぜ、電子が増えても原子は劇的に大きくなるわけではないのか。}
\lecturerresponse{電子を加えることは、原子核に陽子を加えることでもある。増大した核電荷は内側の電子を内へ引き寄せる一方、それらの内側の電子は原子核を遮蔽する。この単純な描像では、外側の価電子はおおよそ1の有効電荷の中を運動することになる。}
\end{classroomqa}

増大した核電荷は内側の電子雲を収縮させる。外側に価電子が1個ある原子では、遠方において原子核と内側の電子がほぼ打ち消し合い、価電子はほぼ単位正電荷のような場の中に残される。したがって原子の大きさは、電子の数に比例して単純に大きくなるわけではない。

\begin{classroomqa}
\lecturetimestamp{01:24:49}
\audiencequestion{隣り合う原子が離れたままであるのは、電気的反発のためか。}
\lecturerresponse{それだけによるのではない。電子雲が重なるとき、パウリ排他が重要な有効反発を与え、多くの原子にとって、それ以上重なることに抵抗する主たる理由はこれである。}
\end{classroomqa}

パウリ反発は排他の有効な帰結であり、余分な力の法則ではない。

\begin{classroomqa}
\lecturetimestamp{01:25:23}
\audiencequestion{排他原理は別個の要請なのか。}
\lecturerresponse{非相対論的量子力学では、それは要請として課される。相対論的量子論では、スピン--統計の関係がそれにより深い基礎を与える。}
\end{classroomqa}

\section{古典力学と量子力学における同種粒子}

その要請を定式化する前に、同種粒子が古典的にどのように数えられるかを考えよう。同じ種類の二つの粒子を一組の箱に入れる。もし粒子が密かに名前---ハリーとサリー---をもっているなら、ハリーが第一の箱に、サリーが第二の箱にある場合は、サリーが第一に、ハリーが第二にある場合とは異なって見える。

古典力学は、そのラベルが検出できないほど薄くても、その描像を許す。各粒子に任意に弱い印を想像し、それでもそれらの軌道を別個のものと見なすことができる。古典統計力学では、通常、同種粒子の置換を新しい配置として数えないことを選ぶ。\(N\) 個の粒子に対して、この慣習はおなじみの因子 \(1/N!\) を生じる。

量子力学はより厳格である。同種粒子は、原理的に追跡できる隠れた個別のラベルをもたない。二つの電子は電子として入れ替え可能である。電子を光子と交換すれば、それらは異なる種であるため物理的に異なる配置を生じるが、二つの電子の席を交換しても、実際にどちらの電子がどちらであったかは明らかにできない。

一粒子状態は波動関数
\[
\psi(x)=\langle x\vert\psi\rangle.
\]
をもつ。二粒子の場合、状態は二つの位置引数の関数であり、
\begin{equation}
\psi(x_{1},x_{2}).
\end{equation}
である。添字は二つの空間方向ではなく、記述における二つのスロットを表す。その絶対値の二乗は結合位置確率密度を与える。

\section{交換演算子と統計の要請}

二つの粒子スロットを入れ替える：
\begin{equation}
\psi(x_{1},x_{2})\longmapsto\psi(x_{2},x_{1}).
\end{equation}
この議論で用いられている記法に従い、交換演算子を\(S\)と呼ぶ。ここで\(S\)は交換を表し、スピンベクトル\(\mathbf S\)ではない。その作用は
\begin{equation}
(S\psi)(x_{1},x_{2})=\psi(x_{2},x_{1}).
\end{equation}
である。二変数の一般関数は対称である必要はないので、この操作は異なる波動関数を生じうる。

二回入れ替えると元の順序に戻る：
\begin{equation}
S^{2}=I.
\end{equation}
交換は可逆であり、内積を保存するので、ユニタリ演算子で表される。もし\(\psi\)が固有値\(\lambda\)をもつ\(S\)の固有ベクトルならば、
\[
S^{2}\psi=\lambda^{2}\psi=\psi,
\]
であり、したがって
\begin{equation}
\lambda=\pm1.
\end{equation}

この結論は二つの交換固有空間を特定するが、任意の二粒子波動関数がすでに\(S\)の固有ベクトルであるとは述べていない。そのような波動関数はその二つの部分に分離でき、\footnote{この分解は、\(S^{2}=I\)によって示唆される二つの固有空間についての編集上の代数的明確化である。記録された議論は二つの固有値を述べ、次に追加の統計原理を導入する。}
\begin{equation}
\psi_{\pm}=\frac12(I\pm S)\psi,
\end{equation}
これは
\begin{equation}
S\psi_{\pm}=\pm\psi_{\pm}.
\end{equation}
を満たす。したがって
\begin{align}
\psi_{+}(x_{2},x_{1})&=\psi_{+}(x_{1},x_{2}),\\
\psi_{-}(x_{2},x_{1})&=-\psi_{-}(x_{1},x_{2}).
\end{align}

ここで、\(S^{2}=I\)だけから導かれる帰結ではなく、追加の原理が現れる。任意の固定された種の同種粒子に対して、許される状態はこれらの交換セクターの一方に完全に属する。すなわち、一対を交換すると常に同じ波動関数が戻るか、常にその負が戻るかのいずれかである。\(+1\)セクターは対称であり、\(-1\)セクターは反対称である。これら二つの可能性を二つの粒子統計として物理的に解釈することは次の段階であり、以下の議論に委ねられる。
